%\title{University of Bristol Thesis Template}
\RequirePackage[l2tabu]{nag}		% Warns for incorrect (obsolete) LaTeX usage
%
%
% File: memoirthesis.tex
% Author: Victor Brena
% Description: Contains the thesis template using memoir class,
% which is mainly based on book class but permits better control of 
% chapter styles for example. This template is an adaptation and 
% modification of Oscar's.
% 
% Memoir is a flexible class for typesetting poetry, fiction, 
% non-fiction and mathematical works as books, reports, articles or
% manuscripts. CTAN repository is found at:
% http://www.ctan.org/tex-archive/macros/latex/contrib/memoir/
%
%
% UoB guidelines for thesis presentation were found at:
% http://www.bris.ac.uk/esu/pg/pgrcop11-12topic.pdf#page=49
%
% UoB guidelines:
%
% The dissertation must be printed on A4 white paper. Paper up to A3 may be used
% for maps, plans, diagrams and illustrative material. Pages (apart from the
% preliminary pages) should normally be double-sided.
%
% Memoir class loads useful packages by default (see manual).
\documentclass[a4paper,11pt,leqno,openbib,oldfontcommands]{memoir} %add 'draft' to turn draft option on (see below)
%
%
% Adding metadata:

\usepackage{datetime}
\usepackage{ifpdf}
\usepackage{acronym}
\usepackage{setspace,lipsum}% http://ctan.org/pkg/{setspace,lipsum}
\usepackage{multirow}
\usepackage[ruled,vlined,algochapter]{algorithm2e}
\usepackage{algorithmic-pt}  % Algoritmos em pt
\usepackage{listing}
\usepackage{listings}
\usepackage{framed}
\usepackage{subfigure}
\usepackage{caption}
   \usepackage[table]{xcolor}% http://ctan.org/pkg/xcolor
  \usepackage{fancybox}
 

\usepackage{amsthm}
 
\theoremstyle{definition}
\newtheorem{definition}{Definici\'on}[section]

\newcommand{\source}[1]{\caption*{Fuente: {#1}} }

\lstset{language=C++,basicstyle=\footnotesize\rmfamily, % Standardschrift
         numbers=left,               % Ort der Zeilennummern
         numberstyle=\tiny,          % Stil der Zeilennummern
         numbersep=5pt,              % Abstand der Nummern zum Text
         tabsize=3,                  % Groesse von Tabs
         extendedchars=true,         %
         breaklines=true,            % Zeilen werden Umgebrochen
         frame=b,
         keywordstyle=[1]\textbf,    % Stil der Keywords
         stringstyle=\color{white}\ttfamily, % Farbe der String
         showspaces=false,           % Leerzeichen anzeigen ?
         showtabs=false,             % Tabs anzeigen ?
         xleftmargin=17pt,
         framexleftmargin=17pt,
         framexrightmargin=5pt,
         framexbottommargin=4pt,
         rulesepcolor=\color{gray},
         columns=fullflexible,
         %extendedchar=true,
         %backgroundcolor=\color{lightgray},
         showstringspaces=false,      % Leerzeichen in Strings anzeigen ?
    emph={synchronizeThreadGroup}, emphstyle={\textbf},
    emph={[2]atomicInc}, emphstyle={[2]\textbf},
    emph={[3]startThreads}, emphstyle={[3]\textbf},
    emph={[4]end}, emphstyle={[4]\textbf},
    emph={[5]algorithm}, emphstyle={[5]\textbf},
    emph={[6]deviceMem},emphstyle={[6]\textbf},
    emph={[7]kernel},emphstyle={[7]\textbf},
    emph={[8]to},emphstyle={[8]\textbf},
    emph={[9]then},emphstyle={[9]\textbf},
    emph={[10]procedure},emphstyle={[10]\textbf},
    emph={[11]dataset},emphstyle={[11]\textbf},
    emph={[12]shared},emphstyle={[12]\textbf},
    emph={[13]Entrada},emphstyle={[13]\textbf},
    emph={[14]Saida},emphstyle={[14]\textbf},
    emph={[15]foreach},emphstyle={[15]\textbf},
    emph={[16]atomicExch},emphstyle={[16]\textbf},
    emph={[17]atomicMin},emphstyle={[17]\textbf},
    emph={[18]shared},emphstyle={[18]\textbf}
 }
 
  \renewcommand*\lstlistingname{Algoritmo}
 \renewcommand*\lstlistlistingname{Lista de Algoritmos}
  
\ifpdf
\pdfinfo{
   /Author (Author's name)
   /Title (PhD Thesis)
   /Keywords (One; Two;Three)
   /CreationDate (D:\pdfdate)
}
\fi
% When draft option is on. 
\ifdraftdoc 
	\usepackage{draftwatermark}				%Sets watermarks up.
	\SetWatermarkScale{0.3}
	\SetWatermarkText{\bf Draft: \today}
\fi
%
% Declare figure/table as a subfloat.
\newsubfloat{figure}
\newsubfloat{table}
% Better page layout for A4 paper, see memoir manual.
\settrimmedsize{297mm}{210mm}{*}
\setlength{\trimtop}{0pt} 
\setlength{\trimedge}{\stockwidth} 
\addtolength{\trimedge}{-\paperwidth} 
\settypeblocksize{634pt}{448.13pt}{*} 
\setulmargins{4cm}{*}{*} 
\setlrmargins{*}{*}{1.5} 
\setmarginnotes{17pt}{51pt}{\onelineskip} 
\setheadfoot{\onelineskip}{2\onelineskip} 
\setheaderspaces{*}{2\onelineskip}{*} 
\checkandfixthelayout
 

\frenchspacing
% Font with math support: New Century Schoolbook
\usepackage{fouriernc}
\usepackage[T1]{fontenc}
%
% UoB guidelines:
%
% Text should be in double or 1.5 line spacing, and font size should be
% chosen to ensure clarity and legibility for the main text and for any
% quotations and footnotes. Margins should allow for eventual hard binding.
%
% Note: This is automatically set by memoir class. Nevertheless \OnehalfSpacing 
% enables double spacing but leaves single spaced for captions for instance. 
\OnehalfSpacing 
%
% Sets numbering division level
\setsecnumdepth{subsection} 
\maxsecnumdepth{subsubsection}
%
% Chapter style (taken and slightly modified from Lars Madsen Memoir Chapter 
% Styles document
\usepackage{calc,soul,fourier}
\makeatletter 
\newlength\dlf@normtxtw 
\setlength\dlf@normtxtw{\textwidth} 
\newsavebox{\feline@chapter} 
\newcommand\feline@chapter@marker[1][4cm]{%
	\sbox\feline@chapter{% 
		\resizebox{!}{#1}{\fboxsep=1pt%
			\colorbox{gray}{\color{white}\thechapter}% 
		}}%
		\rotatebox{90}{% 
			\resizebox{%
				\heightof{\usebox{\feline@chapter}}+\depthof{\usebox{\feline@chapter}}}% 
			{!}{\scshape\so\@chapapp}}\quad%
		\raisebox{\depthof{\usebox{\feline@chapter}}}{\usebox{\feline@chapter}}%
} 
\newcommand\feline@chm[1][4cm]{%
	\sbox\feline@chapter{\feline@chapter@marker[#1]}% 
	\makebox[0pt][c]{% aka \rlap
		\makebox[1cm][r]{\usebox\feline@chapter}%
	}}
\makechapterstyle{daleifmodif}{
	\renewcommand\chapnamefont{\normalfont\Large\scshape\raggedleft\so} 
	\renewcommand\chaptitlefont{\normalfont\Large\bfseries\scshape} 
	\renewcommand\chapternamenum{} \renewcommand\printchaptername{} 
	\renewcommand\printchapternum{\null\hfill\feline@chm[2.5cm]\par} 
	\renewcommand\afterchapternum{\par\vskip\midchapskip} 
	\renewcommand\printchaptertitle[1]{\color{gray}\chaptitlefont\raggedleft ##1\par}
} 
\makeatother 
\chapterstyle{daleifmodif}
%
% UoB guidelines:
%
% The pages should be numbered consecutively at the bottom centre of the
% page.
\makepagestyle{myvf} 
\makeoddfoot{myvf}{}{\thepage}{} 
\makeevenfoot{myvf}{}{\thepage}{} 
\makeheadrule{myvf}{\textwidth}{\normalrulethickness} 
\makeevenhead{myvf}{\small\textsc{\leftmark}}{}{} 
\makeoddhead{myvf}{}{}{\small\textsc{\rightmark}}
\pagestyle{myvf}
%
% Oscar's command (it works):
% Fills blank pages until next odd-numbered page. Used to emulate single-sided
% frontmatter. This will work for title, abstract and declaration. Though the
% contents sections will each start on an odd-numbered page they will
% spill over onto the even-numbered pages if extending beyond one page
% (hopefully, this is ok).
\newcommand{\clearemptydoublepage}{\newpage{\thispagestyle{empty}\cleardoublepage}}
%
%
% Creates indexes for Table of Contents, List of Figures, List of Tables and Index
\makeindex
% \printglossaries below creates a list of abbreviations. \gls and related
% commands are then used throughout the text, so that latex can automatically
% keep track of which abbreviations have already been defined in the text.
%
% The import command enables each chapter tex file to use relative paths when
% accessing supplementary files. For example, to include
% chapters/brewing/images/figure1.png from chapters/brewing/brewing.tex we can
% use
% \includegraphics{images/figure1}
% instead of
% \includegraphics{chapters/brewing/images/figure1}
\usepackage{import}

% Add other packages needed for chapters here. For example:
\usepackage{lipsum}					%Needed to create dummy text
\usepackage{amsfonts} 					%Calls Amer. Math. Soc. (AMS) fonts
\usepackage[centertags]{amsmath}			%Writes maths centred down
\usepackage{stmaryrd}					%New AMS symbols
\usepackage{amssymb}					%Calls AMS symbols
\usepackage{amsthm}					%Calls AMS theorem environment
\usepackage{newlfont}					%Helpful package for fonts and symbols
\usepackage{layouts}					%Layout diagrams
\usepackage{graphicx}					%Calls figure environment
\usepackage{longtable,rotating}			%Long tab environments including rotation. 
\usepackage[utf8]{inputenc}			%Needed to encode non-english characters 
									%directly for mac
\usepackage{colortbl}					%Makes coloured tables
\usepackage{wasysym}					%More math symbols
\usepackage{mathrsfs}					%Even more math symbols
\usepackage{float}						%Helps to place figures, tables, etc. 
\usepackage{verbatim}					%Permits pre-formated text insertion
\usepackage{upgreek }					%Calls other kind of greek alphabet
\usepackage{latexsym}					%Extra symbols
\usepackage[square,numbers,
		     sort&compress]{natbib}		%Calls bibliography commands 
\usepackage{url}						%Supports url commands
% \usepackage{etex}						%eTeXÕs extended support for counters
% \usepackage{fixltx2e}					%Eliminates some in felicities of the 
									%original LaTeX kernel
\usepackage[spanish,english]{babel}		%For languages characters and hyphenation
\usepackage[spanish]{babel}		%For languages characters and hyphenation

\usepackage{color}                    				%Creates coloured text and background
\usepackage[colorlinks=true,
		     allcolors=black]{hyperref}              %Creates hyperlinks in cross references
\usepackage{memhfixc}					%Must be used on memoir document 
									%class after hyperref
\usepackage{enumerate}					%For enumeration counter
\usepackage{footnote}					%For footnotes
\usepackage{microtype}					%Makes pdf look better.
\usepackage{rotfloat}					%For rotating and float environments as tables, 
									%figures, etc. 
\usepackage{alltt}						%LaTeX commands are not disabled in 
									%verbatim-like environment
\usepackage[version=0.96]{pgf}			%PGF/TikZ is a tandem of languages for producing vector graphics from a 
\usepackage{tikz}						%geometric/algebraic description.
\usetikzlibrary{arrows,shapes,snakes,
		       automata,backgrounds,
		       petri,topaths}				%To use diverse features from tikz		
		       
		       
\selectlanguage{spanish}

\addto\captionsenglish{% Replace "english" with the language you use
  \renewcommand{\contentsname}%
    {Índice de Contenidos}%
} 
    
%							
%Reduce widows  (the last line of a paragraph at the start of a page) and orphans 
% (the first line of paragraph at the end of a page)
\widowpenalty=1000
\clubpenalty=1000
%
% New command definitions for my thesis
%
\newcommand{\keywords}[1]{\par\noindent{\small{\bf Keywords:} #1}} %Defines keywords small section
\newcommand{\parcial}[2]{\frac{\partial#1}{\partial#2}}                             %Defines a partial operator
\newcommand{\vectorr}[1]{\mathbf{#1}}                                                        %Defines a bold vector
\newcommand{\vecol}[2]{\left(                                                                         %Defines a column vector
	\begin{array}{c} 
		\displaystyle#1 \\
		\displaystyle#2
	\end{array}\right)}
\newcommand{\mados}[4]{\left(                                                                       %Defines a 2x2 matrix
	\begin{array}{cc}
		\displaystyle#1 &\displaystyle #2 \\
		\displaystyle#3 & \displaystyle#4
	\end{array}\right)}
\newcommand{\pgftextcircled}[1]{                                                                    %Defines encircled text
    \setbox0=\hbox{#1}%
    \dimen0\wd0%
    \divide\dimen0 by 2%
    \begin{tikzpicture}[baseline=(a.base)]%
        \useasboundingbox (-\the\dimen0,0pt) rectangle (\the\dimen0,1pt);
        \node[circle,draw,outer sep=0pt,inner sep=0.1ex] (a) {#1};
    \end{tikzpicture}
}
\newcommand{\range}[1]{\textnormal{range }#1}                                             %Defines range operator
\newcommand{\innerp}[2]{\left\langle#1,#2\right\rangle}                                 %Defines inner product
\newcommand{\prom}[1]{\left\langle#1\right\rangle}                                         %Defines average operator
\newcommand{\tra}[1]{\textnormal{tra} \: #1}                                                       %Defines trace operator
\newcommand{\sign}[1]{\textnormal{sign\,}#1}                                                   %Defines sign operator
\newcommand{\sech}[1]{\textnormal{sech} #1}                                                  %Defines sech
\newcommand{\diag}[1]{\textnormal{diag} #1}                                                    %Defines diag operator
\newcommand{\arcsech}[1]{\textnormal{arcsech} #1}                                       %Defines arcsech
\newcommand{\arctanh}[1]{\textnormal{arctanh} #1}                                         %Defines arctanh
%Change tombstone symbol
\newcommand{\blackged}{\hfill$\blacksquare$}
\newcommand{\whiteged}{\hfill$\square$}
\newcounter{proofcount}
\renewenvironment{proof}[1][\proofname.]{\par
 \ifnum \theproofcount>0 \pushQED{\whiteged} \else \pushQED{\blackged} \fi%
 \refstepcounter{proofcount}
 \normalfont 
 \trivlist
 \item[\hskip\labelsep
       \itshape
   {\bf\em #1}]\ignorespaces
}{%
 \addtocounter{proofcount}{-1}
 \popQED\endtrivlist
}
%
%
% New definition of square root:
% it renames \sqrt as \oldsqrt
\let\oldsqrt\sqrt
% it defines the new \sqrt in terms of the old one
\def\sqrt{\mathpalette\DHLhksqrt}
\def\DHLhksqrt#1#2{%
\setbox0=\hbox{$#1\oldsqrt{#2\,}$}\dimen0=\ht0
\advance\dimen0-0.2\ht0
\setbox2=\hbox{\vrule height\ht0 depth -\dimen0}%
{\box0\lower0.4pt\box2}}
%
% My caption style
\newcommand{\mycaption}[2][\@empty]{
	\captionnamefont{\scshape} 
	\changecaptionwidth
	\captionwidth{0.9\linewidth}
	\captiondelim{.\:} 
	\indentcaption{0.75cm}
	\captionstyle[\centering]{}
	\setlength{\belowcaptionskip}{10pt}
	\ifx \@empty#1 \caption{#2}\else \caption[#1]{#2}
}
%
% My subcaption style
\newcommand{\mysubcaption}[2][\@empty]{
	\subcaptionsize{\small}
	\hangsubcaption
	\subcaptionlabelfont{\rmfamily}
	\sidecapstyle{\raggedright}
	\setlength{\belowcaptionskip}{10pt}
	\ifx \@empty#1 \subcaption{#2}\else \subcaption[#1]{#2}
}
%
%An initial of the very first character of the content
\usepackage{lettrine}
\newcommand{\initial}[1]{%
	\lettrine[lines=3,lhang=0.33,nindent=0em]{
		\color{gray}
     		{\textsc{#1}}}{}}
%
% Theorem styles used in my thesis
%
\theoremstyle{plain}
\newtheorem{theo}{Theorem}[chapter]
\theoremstyle{plain}
\newtheorem{prop}{Proposition}[chapter]
\theoremstyle{plain}
\theoremstyle{definition}
\newtheorem{dfn}{Definition}[chapter]
\theoremstyle{plain}
\newtheorem{lema}{Lemma}[chapter]
\theoremstyle{plain}
\newtheorem{cor}{Corollary}[chapter]
\theoremstyle{plain}
\newtheorem{resu}{Result}[chapter]
%
% Hyphenation for some words
%
\hyphenation{res-pec-tively}
\hyphenation{mono-ti-ca-lly}
\hyphenation{hypo-the-sis}
\hyphenation{para-me-ters}
\hyphenation{sol-va-bi-li-ty}
%
%

 

 \newcommand{\Nhashfunctions}{m}
 \newcommand{\TopKMotifs}{\textit{TopK-Motifs} }
 \newcommand{\kNN}{kNN }
 \newcommand{\AllkNN}{AllkNN }
 \newcommand{\CUDAAllkNN}{\emph{CUDA-AllkNN()} }
 \newcommand{\CUDATopkMotifs}{\emph{CUDA-TopKMotifs()} }

\acrodef{CBIR}{Content-based image retrieval}
\acrodef{MAM}{Metric Access Method}
\acrodef{AM}{Access Method}
\acrodef{SAM}{Spatial Access Method}
\acrodef{MST} {Minimum Spanning Tree}
\acrodef{IR} {Information Retrieval}
\acrodef{NDC} {Number of Distance Calculations}
\acrodef{LAESA} {Linear Approximating Eliminating Search Algorithm}
\acrodef{AESA} {Approximating Eliminating Search Algorithm}
\acrodef{RNG} {Relative Neighborhood Graph}
\acrodef{DT} {Delaunay Triangulation}
\acrodef{GG} {Gabriel Graph}
\acrodef{UG} {Urquhart Graph}
\acrodef{SAT} {Spacial Approximation Tree}
\acrodef{RNDF} {Relative Neighborhood Density Factor}
\acrodef{NN} {Nearest Neighbor}
\acrodef{NAG} {Neighborhood Approximation Graph}
\acrodef{HRG} {Hyperspherical Region Graph}
\acrodef{NA-Graph} {Neighborhood Approximation Graph}
\acrodef{DA} {Disk Access}
\acrodef{DTW} {Dynamic Time Warping}
\acrodef{LSH} {Locality Sensitive Hashing}
\acrodef{CNN} {Red Neuronal  Convolutiva}
\acrodef{CNNs} {Redes Neuronales  Convolutivas}

%\newcommand{ \Correction }[ 1 ] {\textcolor{black}{#1}}
%\newcommand{ \NewMaterial }[ 1 ] {\textcolor{black}{#1}}
%\newcommand{ \ConnectDots }[ 1 ] {\textcolor{black}{#1}}
%\newcommand{ \NewMaterialGPU }[ 1 ] {\textcolor{black}{#1}}
%\newcommand{ \Error }[ 1 ] {\textcolor{black}{#1}}


\acrodef{CBIR} {Content-based Image Retrieval}
\acrodef{MAM}  {Métodos de Acceso Métrico}
\acrodef{SAM}  {Métodos de Acceso Espacial}
\acrodef{MAE}  {Métodos de Acesso Espaciais}

\acrodef{DTW}  {Dynamic Time Warping}
\acrodef{LSH}  {Locality Sensitive Hashing}

\begin{document}
\selectlanguage{spanish}

 % UoB guidlines:
%
% Preliminary pages
% 
% The five preliminary pages must be the Title Page, Abstract, Dedication
% and Acknowledgements, Author's Declaration and Table of Contents.
% These should be single-sided.
% 
% Table of contents, list of tables and illustrative material
% 
% The table of contents must list, with page numbers, all chapters,
 % sections and subsections, the list of references, bibliography, list of
% abbreviations and appendices. The list of tables and illustrations
% should follow the table of contents, listing with page numbers the
% tables, photographs, diagrams, etc., in the order in which they appear
% in the text.
% 
 
\renewcommand{\listtablename}{Índice de tablas}

\frontmatter
\pagenumbering{roman}
%
%
% File: Title.tex
% Author: V?ctor Bre?a-Medina
% Description: Contains the title page
%
% UoB guidelines:


% 
% At the top of the title page, within the margins, the dissertation should give the title and, if 
% necessary, sub-title and volume number. If the dissertation is in a language other than English, the 
% title must be given in that language and in English. The full name of the author should be in the 
% centre of the page. At the bottom centre should be the words ?A dissertation submitted to the 
% University of Bristol in accordance with the requirements for award of the degree of ? in the 
% Faculty of ...?, with the name of the school and month and year of submission. The word count of 
% the dissertation (text only) should be entered at the bottom right-hand side of the page.
%
%
\begin{titlingpage}
\begin{SingleSpace}
\calccentering{\unitlength} 
\begin{adjustwidth*}{\unitlength}{-\unitlength}

\begin{center}
\fontsize{20pt}{20pt} \selectfont UNIVERSIDAD NACIONAL DE SAN AGUSTÍN 
ESCUELA DE POSTGRADO

\vspace*{1cm}
\begin{center}
{\Large UNIDAD DE POSTGRADO DE LA FACULTAD DE INGENIERÍA 
DE PRODUCCIÓN Y SERVICIOS}
\end{center}
\vspace*{1cm} 

\begin{center}
      
    
      \includegraphics[scale=0.77]{imagenes/unsaescudo.jpg}\\ 
\end{center}
      

\vspace*{1cm}
\begin{center}
{\Large MÉTODOS SEMÁNTICOS PARA LA RECUPERACIÓN DE INFORMACIÓN EN GRANDES VOLÚMENES DE DATOS: UNA ARQUITECTURA ESCALABLE Y EFICIENTE }

\end{center} 


\vspace*{1.5cm}
\begin{flushright}
 
\small{Tesis presentado por el Magister\\
\vspace*{0.15cm}

\textbf{Alexander Victor Ocsa Mamani}\\
\vspace*{0.2cm}

Para optar el Grado de Doctor en Ciencia de la Computación\\
\vspace*{0.15cm}
Asesor: Dr. Cristian Lopez del Alamo}
 
\end{flushright}

\vspace*{2.5cm}
\begin{center}
\small
\textbf{Arequipa - Perú \\2017}
\end{center}
\cleardoublepage

\end{center}
\end{adjustwidth*}
\end{SingleSpace}
\end{titlingpage}



\newpage





\vfill

\vspace*{5cm}

\begin{center}
\begin{minipage}[c]{12cm}
\begin{center}
\hrulefill\\
\vspace{.5cm} {\large MÉTODOS SEMÁNTICOS PARA LA RECUPERACIÓN DE INFORMACIÓN EN GRANDES VOLÚMENES DE DATOS: UNA ARQUITECTURA ESCALABLE Y EFICIENTE}\\
\vspace{1.3cm}
\small\textbf{\it Alexander Victor Ocsa Mamani}\\
\vspace{.5cm}
\hrulefill\\
\end{center}
\end{minipage}
\end{center}

\vfill

 



 
\clearemptydoublepage
%
%
% File: abstract.tex
% Author: V?ctor Bre?a-Medina
% Description: Contains the text for thesis abstract
%
% UoB guidelines:
%
% Each copy must include an abstract or summary of the dissertation in not
% more than 300 words, on one side of A4, which should be single-spaced in a
% font size in the range 10 to 12. If the dissertation is in a language other
% than English, an abstract in that language and an abstract in English must
% be included.

\chapter*{Resumen}
\begin{SingleSpace}
\initial{L}a creciente disponibilidad de datos en diferente dominio de aplicación ha motivado el desarrollo de técnicas de recuperación y descubrimiento de conocimiento en grandes volúmenes de datos.   Recientes trabajos muestran que tanto las técnicas de aprendizaje profundo como  nuevos métodos de búsqueda aproximada en dominio de datos complejos son campos de investigación importantes, donde tanto la eficiencia como la  escalabilidad de los algoritmos son factores críticos. Para resolver el problema de escalabilidad se han propuesto muchos enfoques. En problemas de gran escala con datos en altas dimensiones, una solución de búsqueda aproximada con un análisis teórico solido se muestra más adecuado que una solución exacta con un modelo teórico débil.    Algoritmos de búsqueda aproximada  basados en \textit{hashing} son propuestos para consultar en conjuntos de datos  alta dimensiones debido a su velocidad de recuperación y bajo costo de almacenamiento.  Por otro lado, en problemas donde se tiene grandes volúmenes de datos etiquetados las técnicas de aprendizaje profundo, como las redes convolucionales, se muestran más adecuadas conforme el número de ejemplos por clases crece.

Estudios recientes, promueven el uso de la \acf{CNN} con técnicas de  \textit{hashing} para mejorar la precisión de la búsqueda de los k-vecinos más cercanos - KNN.  Sin embargo, aun hay retos que resolver para encontrar una solución práctica y eficiente para indexar características  CNN, tales como la necesidad de un proceso de entrenamiento intenso para lograr resultados de consulta precisos y la dependencia crítica de los parámetros.   Con el fin de superar estos problemas, se propone un nuevo método de búsqueda por similitud, \textit{Deep frActal based  Hashing (DAsH)}, para calcular los mejores valores de parámetros  para una proyección óptima en un subespacio, explorando las correlaciones entre los atributos de las características  CNN usando la teoría fractal. Además, inspirado por recientes avances  en redes CNN, utilizamos no solo activaciones de capas inferiores que son más generales, sino también el conocimiento previo de los datos semánticos sobre la última capa CNN para mejorar la precisión de la búsqueda.  Así, nuestro método produce una mejor representación del espacio de datos con un coste computacional menor para una mejor precisión. Esta mejora  significativa en velocidad y precisión nos permite evaluar este esquema en conjuntos de datos reales y sintéticos.

 
\paragraph{Palabras claves:}
Recuperación de Información, Búsqueda kNN, LSH, Deep Learning,  CNN, Dimensión Fractal.


\end{SingleSpace}
\clearpage



\chapter*{Abstract}
\begin{SingleSpace}
\initial{T}he increasing availability of data in different application domains has led to the development of techniques for recovering and discovering knowledge in large volumes of data.Recent works shows that both deep learning techniques and new methods of approximate searching in complex data domain are important research fields, where both the efficiency and scalability of algorithms are critical factors. Many approaches have been proposed to solve the problem of scalability. In large scale problems with high dimensional data, an approximate search solution with a solid theoretical analysis is more appropriate than an exact solution with a weak theoretical model.    Approximate search algorithms based on |textittit{hashing} are proposed to query in large datasets due to their speed of retrieval and low storage cost.  On the other hand, in problems where there are large volumes of labelled data, deep learning techniques, such as convolutional networks, are more appropriate as the number of examples per class grows.

Recent studies, promote the use of  \ac{CNN} with hashing techniques to improve the search accuracy for nearest neighbor search. However, there are challenges to solve in order to find a practical and efficient solution to index CNN features, such as the need for heavy training process to achieve accurate query results and the critical dependency on data-parameters.   Aiming to overcome these issues, we propose a new method for scalable similarity search, i.e., Deep frActal based  Hashing (DAsH), by computing the best data-parameters values for optimal sub-space projection  exploring the correlations among CNN features attributes using fractal theory. Moreover, inspired by recent advances in CNNs, we use not only activations of lower layers which are more general-purpose but also previous knowledge of the semantic data on the latest CNN layer to improve the search accuracy.  Thus, our method produces a better representation of the data space with a less computational cost for a better accuracy.  This significant gain in speed and accuracy allows us to evaluate the framework on a  large, realistic, and challenging set of datasets.  
 
\paragraph{Keywords:}
Information retrieval, kNN search, LSH, Deep Learning, CNN, Fractal dimension. 

\end{SingleSpace}
\clearpage
\clearemptydoublepage
%
%
% file: dedication.tex
% author: V?ctor Bre?a-Medina
% description: Contains the text for thesis dedication
%

\chapter*{Agradecimientos}
\begin{SingleSpace}
\initial{A} mi orientador, Profesor. Dr. Cristian Lopez del Alamo, por estar siempre presente y atento. Por todos los años de orientación, de los trabajos de investigación que concluye esta tesis de doctorado. Por los desafíos que me proporcionó, y principalmente por su confianza. A mis padres y a toda mi familia por los años de educación que me proporcionaron. A mi  novia  Karla, por todo el apoyo que siempre me ha brindado. A Concytec y Cienciactiva por todo el apoyo y darnos la confianza de realizar investigación.  A todos los profesores que nos proporcionaron los conocimientos necesarios para realizar investigación, y finalmente a todas las personas que directa o indirectamente contribuyeron para que llegase hasta aquí.
\end{SingleSpace}
\clearpage
\clearemptydoublepage
%
%%
% File: declaration.tex
% Author: V?ctor Bre?a-Medina
% Description: Contains the declaration page
%
% UoB guidelines:
%
% Author's declaration
%
% I declare that the work in this dissertation was carried out in accordance
% with the requirements of the University's Regulations and Code of Practice
% for Research Degree Programmes and that it has not been submitted for any
% other academic award. Except where indicated by specific reference in the
% text, the work is the candidate's own work. Work done in collaboration with,
% or with the assistance of, others, is indicated as such. Any views expressed
% in the dissertation are those of the author.
%
% SIGNED: .............................................................
% DATE:..........................
%
\chapter*{Author's declaration}
\begin{SingleSpace}
\begin{quote}
\initial{I} declare that the work in this dissertation was carried out in accordance with the requirements of  the University's Regulations and Code of Practice for Research Degree Programmes and that it  has not been submitted for any other academic award. Except where indicated by specific  reference in the text, the work is the candidate's own work. Work done in collaboration with, or with the assistance of, others, is indicated as such. Any views expressed in the dissertation are those of the author.

\vspace{1.5cm}
\noindent
\hspace{-0.75cm}\textsc{SIGNED: .................................................... DATE: ..........................................}
\end{quote}
\end{SingleSpace}
\clearpage
%\clearemptydoublepage
%
\renewcommand{\contentsname}{Índice General}
\maxtocdepth{subsection}
\tableofcontents*
\addtocontents{toc}{\par\nobreak \mbox{}\hfill{\bf Página}\par\nobreak}
\clearemptydoublepage
%
\listoftables
\addtocontents{lot}{\par\nobreak\textbf{{\scshape Tabla} \hfill Página}\par\nobreak}
\clearemptydoublepage
%
\listoffigures
\addtocontents{lof}{\par\nobreak\textbf{{\scshape Figura} \hfill Página}\par\nobreak}
\clearemptydoublepage


\addcontentsline{toc}{chapter}{Siglas}
\chapter*{Siglas}

 \begin{tabular}{ l l }
     
KDD & \textit{Knowledge Discovery in Databases}.\\
LSH & \textit{Locality Sensitive Hashing}.\\
MA & Método de Acesso.\\
MAM & Método de Acesso Métrico.\\
MAE & Método de Acesso Espacial.\\
NCD & Número de Calculos de Distancia.\\
GEMINI &  \textit{GEneric Multimedia INdexIng}.\\
ML & \textit{Machine Learning}.\\
CNN & \textit{Convolutional Neural Network}.\\

\end{tabular}
\cleardoublepage


\addcontentsline{toc}{chapter}{Lista de símbolos}
\chapter*{Lista de símbolos}

 \begin{tabular}{ l l l l } 
  %\hline
     $S$ & conjunto de dados\\
    $N$ & número de objetos en el  conjunto de dados\\
     $q$ & objeto de consulta\\
    $r$ & radio de consulta\\
    $k$ & n\'umero de objetos m\'as pr\'oximos en una consulta kNN\\
    $\varepsilon$ & error relativo en la busca por similaridade aproximada\\
    $(1 + \varepsilon)$ & o factor de aproximaci\'on \\
    
    $Rq(q, r)$  & consulta por rango \\
    $kNN(q, k)$ & consulta aos k-vecinos mais pr\'oximos\\
     
    $(1+\epsilon)-Rq(q,r)$ & consulta por rango aproximada \\
    $(1+\epsilon)-kNN(q,k)$ & consulta aproximada a los k-vecinos m\'as pr\'oximos \\
 
    $h()$ & funci\'on \textit{hash} \\
    $H = \{h_1(),...,h_\Nhashfunctions()\}$ & un conjunto de funciones \textit{hash} de tama\~no $\Nhashfunctions$ \\
    $L$ & número de tablas \textit{hash}\\
    $\Nhashfunctions$ & n\'umero de funciones  \textit{hash} (dimensi\'on de la proyecci\'on \textit{hash})\\
     $\tau$ & número de \textit{buckets} o $probes$ a ser explorados \\
    $X = (x_1, x_2, ...,x_n)$ & vector característica \\
    $E$ &  dimensión embebida de un conjunto de datos \\  
  	$\mathfrak{D}$ & dimension intrínseca (dimensión fractal)
 
\end{tabular}
\cleardoublepage


%
%
% The bulk of the document is delegated to these chapter files in
% subdirectories.
\mainmatter
%
\import{}{chapter1.tex}
\import{}{chapter2.tex}
\import{}{chapter3.tex}
\import{}{chapter4.tex}
\import{}{chapter5.tex}
\import{}{chapter6.tex}
\import{}{chapter7.tex}

\clearemptydoublepage
%
%
% And the appendix goes here
 %
% Apparently the guidelines don't say anything about citations or
% bibliography styles so I guess we can use anything.
\backmatter
 \refstepcounter{chapter}
 \bibliographystyle{siam}

\bibliography{biblio}
\clearemptydoublepage
%
% Add index
%\printindex
%    

\end{document}